% Dokumentanklasse: a4paper, 14pt
% Beschreibung:     Dokumentenformat
% Option:           extraarticle - ?
\documentclass[a4paper,14 pt]{extarticle}

% Paket:            a4paper
% Beschreibung:     A4 Seitenabstände
% Option:           geometry
\usepackage[a4paper]{geometry}

% Paket:            ansmath
% Beschreibung:     Zum darstellen von mathematischen Formeln
\usepackage{amsmath}

% Paket:            ngerman
% Beschreibung:     Deutsche Rechtschreibung
% Option:           babel - Sibentrennung
%\usepackage{ngerman}
\usepackage[ngerman]{babel}

% Paket:            utf8
% Beschreibung:     Stellt Umlaute richtig dar
% Option:           inputenc - Erlaubt die Darstellung der gleichen Zeichen (Character) wie sie in strin überliefert werden
\usepackage[utf8]{inputenc}

% Paket:            makeindex
% Beschreibung:     Ermöglicht das Indexieren von Wörter und den Befehl \printindex um den Index auszugeben
\usepackage{makeidx}
\makeindex

% Paket:            natbib
% Beschreibung:     Für Zitate
% Option:           round - ?
%\usepackage[round]{natbib}

% Paket:            fancyhdr
% Beschreibung:     Ermöglich ein generelles Seitenlayout ein zu stellen mit Kopf und Fußzeile.
\usepackage{fancyhdr}

% Paket:            graphicx
% Beschreibung:     Einbinden von Bildern
% Option:
\usepackage{graphicx}

% Paket:            enumitem
% Beschreibung:     Zeilenabstände bei Aufzählungen definieren
% Option:
\usepackage{enumitem}

% Paket:            float
% Beschreibung:     Zum Ausrichten von Tabellen und Spalten bzw deren zentrierung
% Option:
% Restriktion:      Muss von Paket hyperref geladen werden. Ansonsten funktioniert das Paket nicht.
\usepackage{float}

% Paket:            appendix
% Beschreibung:     Das Paket dient dazu, ausschließlich das Thema einer Überschrift in das Inhaltsverzeichnis zu überführen
% Option:           appendix - Überführt die Überschriften des Anhangs richtig ins das Inhaltsverzeichnis
\usepackage[titletoc]{appendix}

% Paket:            setspace
% Beschreibung:     Setz über die optionen den Zeilenabstand
% Optionen:         Möglicher Zeilenabstand
%                   singlespacing = 1,0
%                   onehalfspacing = 1,5
%                   doublespacing = 2,0
% Restriktion:      Muss von Paket hyperref geladen werden. Ansonsten funktioniert das Paket nicht.
\usepackage[onehalfspacing]{setspace}

% Packet:           Hyperref
% Beschreibung:     Importiert hyperref um Querverweise mittels \hyperref zu erzeugen.
\usepackage[]{hyperref}
\hypersetup{
	pdftitle={Reliable and scalable messaging in Microservices},
	pdfauthor={Markus Pesch},
	pdfsubject={Bachelorthesis}
}

% Packet:           Minted
% Beschreibung:     Dient zum highlining von Quellcode wie beispielsweise Java, Bash oder Python.
\usepackage{minted}
\usemintedstyle{emacs}

% Packet:           tabularx
% Beschreibung:     Werden Tabellen mit diesem Paket erstellt, ist es möglich Zeilenumbrüche in einer Zelle zu erzeugen
\usepackage{tabularx}



% Paket:            biblatex
\usepackage[
style=authoryear-icomp,    % Zitierstil
isbn=false,                % ISBN nicht anzeigen, gleiches geht mit nahezu allen anderen Feldern
pagetracker=true,          % ebd. bei wiederholten Angaben (false=ausgeschaltet, page=Seite, spread=Doppelseite, true=automatisch)
maxbibnames=50,            % maximale Namen, die im Literaturverzeichnis angezeigt werden (ich wollte alle)
maxcitenames=3,            % maximale Namen, die im Text angezeigt werden, ab 4 wird u.a. nach den ersten Autor angezeigt
autocite=inline,           % regelt Aussehen für \autocite (inline=\parancite)
block=space,               % kleiner horizontaler Platz zwischen den Feldern
backref=true,              % Seiten anzeigen, auf denen die Referenz vorkommt
backrefstyle=three+,       % fasst Seiten zusammen, z.B. S. 2f, 6ff, 7-10
date=short,                % Datumsformat
backend=biber
]{biblatex}
\setlength{\bibitemsep}{1em}     % Abstand zwischen den Literaturangaben
\setlength{\bibhang}{2em}        % Einzug nach jeweils erster Zeile

\addbibresource{literatur//bibliothek.bib}


% Packet:           glossaries
% Beschreibung:     Glossar einbinden und Glossarbefehle bereitstellen
% Option:           Gebe Glossar auch als section im Inhaltsverzeichnis aus
\usepackage[toc,section=section]{glossaries}
\makeglossaries
% Glossareinträge
\newglossaryentry{gls:2-trier-architektur}
{
    name={2-Trier-Architektur},
    description={Eine 2-Tier-Architektur beschreibt, dass ein Client direkt mit dem Server kommuniziert. Der Server stellt eine Applikationsschicht und eine Datenbankschicht bereit. Ein Beispiel für eine 2-Tier-Architektur wären PC-Spiele bei denen die Spieler im Multiplayer spielen. Die Kommunikation jedes Spielers als Client erfolgt über einen Server mit dem sich alle Teilnehmer verbinden},
    plural={2-Trier-Architekturen}
}


% Acronyme
\newacronym{api}{API}{Application Programming Interface}
\newacronym{cli}{CLI}{Commando Line Interface}


% Start des Dokuments
\begin{document}

  % Importiere weitere .tex Dokumente
  \begin{titlepage}
  \begin{center}
    \begin{large}
      Tutorium SS18
    \end{large}
    
    \begin{huge}
      \begin{singlespace}
            \textbf{Datawarehouse}
      \end{singlespace}
    \end{huge}

    \vspace{0.5cm}

    \begin{figure}[h]
      \centering
      \includegraphics[width=0.85\textwidth]{img//logo.png}
      \label{img:fh-trier-logo}
    \end{figure}

    \vspace{2cm}
    
    Latex Sourcecode on \input{./git-url.tmp} \\
    based on git commit \input{git-id.tmp} from \input{git-date.tmp}
    
    
  \end{center}
\end{titlepage}

  \pagebreak

  % Pagestyle
  % Setze das Seitenlayout auf leer um Fuß und Kopfzeilen zu unterdrücken
  \pagestyle{empty}

  % Importiere weitere .tex Dokumente
  %\section*{Vorwort}
Dieses Dokument soll den Leser den Umgang mit Git näher bringen. Dazu wird Ihm der Umgang mit den gängigen Git-Befehlen näher gebracht und der Workflow von üblichen Online-Repositories erläutert. 

Als Beispiel wird in diesem Dokument ein Workflow eingerichtet, die dem Entwickler ermöglicht Quellcodeverbesserungen an diesen Dokument durch Merge Requests vorzunehmen.

\newpage


  \tableofcontents

  % Pagestyle
  % Setze das Seitenlayout auf fancyhdr um Fuß- und Kopfzeilen zu setzen
  \pagestyle{fancy}

  % Löscht alle Kopf- und Fußzeilen des pagestyles fancyhdr
  \fancyhf{}

  % Fuß- und Kopfzeile des Paketes fancyhdr
  % [L] - Linkeseite      [O] - Ungerade Seitenzahlen         [LE,LO] - Linkeseite, Gerade- und Ungerade Seitenanzahlen
  % [C] - Mitte           [E] - Gerade Seitenanzahlen         [CE]    - Seitenmitte, nur gerade Seitenanzahlen
  % [R] - Rechteseite                                         [RO]    - Rechteseite, nur ungerade Seitenanzahlen
  % \fancyhead    Kopfzeile
  % \fancyfoot    Fußzeile
  \fancyhead[L]{\rightmark}
  \fancyhead[R]{Seite \thepage}

  % Pixelstärke der Kopfzeilenlinie
  \renewcommand{\headrulewidth}{1pt}

  % Setze die Seitenbeginn zurück
  \setcounter{page}{1}

  % Importiere weitere .tex Dokumente
  \section{Einführung}
\label{sec:einfuehrung}
Dieses Kapitel bietet eine Einführung in die Grundbegriffe von Git und deren generelle Grundkonfiguration.

\subsection{Grundbegriffe}
\label{sec:einfuehrung.grundbegriffe}
Um dem Umgang mit Git zu erläutern ist es notwendig, einige Fachbegriffe zu verstehen.

\myparagraph{Versionskontrollsystem}
Ein Versionskontrollsystem dient zur Verwaltung und Versionierung von Software. Bekannte Projekte die unter dem Versionskontrollsystem Git stehen sind beispielsweise der Linux-Kernel und Git selbst. Beide Projekte stehen unter der Lizenz GPL. Der Quellcode kann mittels git bezogen werden. Dazu später mehr.

Bei Versionskontrollsystemen wird unterschieden zwischen zentralen und dezentralen Systemen. Git ist ein dezentrales System. Der Entwickler ist nicht abhängig von dem Server auf dem die Versionshierarchie gespeichert ist. Der Entwickler ist in der Lage mit Git die komplette Versionshierarchie auf sein eigenes System zu klonen und auf dieser Basis zu arbeiten.

\myparagraph{Repository}
Git  verwaltet von jeder Datei unterschiedliche Zustände bzw. Versionen. Diese werden in dem Repository gespeichert. Mithilfe des Repositories ist es möglich, auf jeden Zustand einer Datei zurück zu springen.

\myparagraph{Working Tree}
Alle Modifikationen an Dateien bzw. dem Quellcode werden im Working Tree vorgenommen. Andere Bezeichnungen hierfür sind auch Working Directory oder Workspace. Wobei bei vielen IDEs im Heimatverzeichnis das Verzeichnis workspace erstellt wird um Projekte zu speichern und unter Versionskontrolle zu setzen.

\myparagraph{Commit}
Ein Commit nimmt jede Veränderung, Modifikation oder Erstellung an Dateien im Working Tree auf. Zusätzlich speichert der Commit noch weitere Informationen wie den Benutzernamen und dessen E-Mail mit Datum und optional gpg-Signatur ab.

\myparagraph{HEAD}
HEAD bildet eine Referenz ab, in welchem Zustand der Entwickler sein Working Tree vor findet.

\myparagraph{Branch}
Jede Software kann in ihrem Verlauf mehrere Entwicklungsverzweigungen haben um beispielsweise in jeder Verzweigung einzelne Teilaufgaben wie Patches oder neuen Features von Software zu entwickeln. Diese Verzweigungen werden Branches genannt. Jede Verzweigung kann später durch einen Merge oder Rebase wieder zurück geführt werden.

\subsection{Grundeinrichtung}
\label{sec:einfuehrung.git}
Wie bereits in in Kapitel \ref{sec:einfuehrung.grundbegriffe} erklärt, speichert Git bei jedem Commit den Benutzernamen und E-Mail ab um nachzuvollziehen von wem die Änderung des Quellcodes stammt. Bevor nun dieses Dokument, das als \LaTeX{} Projekt angelegt wurde, durch Git verwaltet wird, ist es Sinnvoll den Benutzernamen und die E-Mail Adresse anzupassen \footcite{git-1.6-your-identity}.

\begin{minted}[framesep=2mm, fontsize=\small]{bash}
$ git config --global user.name "Hugo McKinnock"
$ git config --global user.email "hugo.mckinnock@example.com"
\end{minted}

\subsection{Signieren von Commits}
\label{sec:einfuehrung.sign-commits}
Da jede Person in der Lage ist den Quellcode boshaft zu ändern unter beliebigen Namen, bietet Git die Möglichkeit an jeden Commit mit einem GnuPG Schlüssel zu signieren.

\begin{info-popup}
  Manche Projekte setzen signierte Commits voraus, wie beispielsweise das Projekt um den Linux Kernel. Nicht signierte Commits werden durch die Hauptverantwortlichen nicht übernommen. Dazu später mehr.
\end{info-popup}

Ist GnuPG auf dem eigenen System nicht vorhanden, kann dies nachträglich installiert werden. Debian basierende Derivate  bieten GnuPG durch die offiziellen Paketquellen an. Dazu reicht die Installation über den Terminalbefehl \textit{sudo apt-get update; sudo apt-get install gpg2 --yes}.

Wird ein Windows-System genutzt, müssen die GnuPG Binaries manuell heruntergeladen und installiert werden. Die Binaries können unter \href{https://www.gpg4win.de}{gpg4win.de} bezogen werden.

Der GnuPG Schlüssel wird anschließend über den Aufruf von gpg im Terminal (Linux Derivate) oder der Git-Bash (Windows) erzeugt. Im Beispiel wird ein RSA 2048 Bit Schlüssel erzeugt.

\begin{verbatim}
$ gpg --gen-key
gpg (GnuPG) 2.1.11; Copyright (C) 2016 Free Software Foundation, Inc.
This is free software: you are free to change and redistribute it.
There is NO WARRANTY, to the extent permitted by law.

Hinweis: "gpg --full-gen-key" ruft den erweiterten Dialog auf.

GnuPG erstellt eine User-ID, um Ihren Schlüssel identifizierbar zu machen.

Ihr Name: Hugo McKinnock
Email-Adresse: hugo.mckinnock@example.com
Sie haben diese User-ID gewählt:
"Hugo McKinnock <hugo.mckinnock@example.com>"

Ändern: (N)ame, (E)-Mail oder (F)ertig/(A)bbrechen? F
gpg: Schlüssel F3FFC9EC ist als ultimativ vertrauenswürdig gekennzeichnet
\end{verbatim}

Wurde der GnuPG Schlüssel erzeugt, kann dieser nun hinterlegt werden um Commits im Hintergrund automatisch zu signieren. Dazu muss zuerst die Schlüssel-ID ausfindig gemacht werden und anschließend Git mitgeteilt werden.

\begin{verbatim}
$ gpg --list-secret-keys --keyid-format 0xSHORT
/home/hugo/.config/gnupg/pubring.gpg
--------------------------------------
sec   rsa2048/F3FFC9EC 2018-03-15 [SC]
uid      [  ultimativ] Hugo McKinnock <hugo.mckinnock@example.com>
ssb   rsa2048/4A1B6E12 2018-03-15 [E]

$ git config --global user.signingkey "0x85ED78DE"
$ git config --global commit.gpgSign true
\end{verbatim}

\begin{info-popup}
  Sollte das Betriebssystem gpg nicht finden, wird die Binary nicht durch die Umgebungsvariable \textit{PATH} aufgelöst. Entweder wird der Pfad zur gpg Binary der Umgebungsvariable \textit{PATH} hinzugefügt oder der absolute Pfad mit Angabe der gpg Binary Git durch folgenden Befehl mitgeteilt - \textit{git config --global gpg.program $ < $Pfad$ > $}
\end{info-popup}

  % Glossar
  \printglossaries
  \newpage

  % Abbildungsverzeichnis
  \listoffigures
  \newpage

  % Literaturverzeichnis
  \printbibliography

  % Erklärung
  %\include{erklaerung//erklaerung}

\end{document}

